\section{Engineering Discussion and Numerical Reliability}
\label{sec:engineering-discussion}

\subsection{Analytical-to-Full-Wave Adjustment}

The final patch width is effectively unchanged from the analytical starting value,
while the final patch length is 1.59~mm shorter than the analytical estimate. This is a
1.45\% reduction. Such adjustment is consistent with the role of the analytical model
as a first-order starting point: the full-wave solution incorporates the actual inset,
feed transition, finite substrate and ground, port geometry, and numerical boundary
configuration.

\subsection{Impedance Matching}

The inset location controls the resistance sampled along the patch. The analytical
38.13~mm estimate and the final 38.2~mm CST value agree closely, indicating that the
cosine-squared resistance model provided a useful initial feed location. The final
$S_{11}$ marker demonstrates a well-defined simulated match near the design frequency.
It does not establish tolerance robustness or measured matching after fabrication.

\subsection{Radiation Behavior}

The 3D and principal-plane plots show a broadside pattern consistent with the expected
radiation from the fringing fields at the patch edges. The two principal planes have
different 3-dB widths, which is expected because the rectangular aperture dimensions
and field distributions are not identical in the two planes.

\subsection{Numerical Qualifications}

The following limitations materially affect interpretation:

\begin{enumerate}
    \item The public archive does not contain the native CST model, raw $S$-parameter
    export, raw far-field data, mesh statistics, or solver log.

    \item No mesh-convergence or boundary-distance convergence study is preserved.
    The final resonance and far-field values therefore represent the documented model,
    not a formally convergence-qualified solution.

    \item The field and far-field monitors are at 0.915~GHz, whereas the matching
    minimum occurs at 0.9186~GHz. The two frequencies differ by only 0.393\%, but the
    distinction is retained explicitly.

    \item PEC metallization omits conductor loss. The reported directivity is not gain,
    and no finite-conductivity sensitivity study is available.

    \item The slightly positive displayed radiation efficiency in dB is not physically
    admissible as an exact efficiency result and is treated as a small numerical-display
    artifact. Efficiency values are therefore not promoted as verified KPIs.

    \item No fabrication, connector transition, manufacturing tolerance, material
    tolerance, VNA calibration, or antenna-range validation is included.
\end{enumerate}

These limitations do not invalidate the documented design workflow, resonance result,
or qualitative broadside pattern. They define the boundary between a successful
simulation study and a measurement-qualified antenna design.
